IRS-assisted secure transmission and IRS-NOMA are active topics, but most IEEE studies still isolate one or two practical constraints at a time. Hardware-aware IRS beamforming was quantified by Shen \emph{et al.}~\cite{shen2021_tcom}, establishing a key impairment baseline. IRS-NOMA joint beamforming was formulated by Mu \emph{et al.}~\cite{mu2020_twc} and low-complexity IRS-NOMA transmission rules were proposed by Ding and Poor~\cite{ding2020_lcomm}. On the security side, Wang \emph{et al.}~\cite{wang2022_twc} introduced IRS-assisted jamming for secure NOMA, while Zhang \emph{et al.}~\cite{zhang2022_tcom} and Jia \emph{et al.}~\cite{jia2023_lcomm} focused on secrecy-centric IRS-NOMA configurations under practical CSI assumptions.

Table~\ref{tab:related_work} summarizes the coverage of these IEEE works versus the present manuscript. The key gap we address is the combined treatment of secrecy-aware optimization, imperfect CSI, and hardware loss within a single IRS-NOMA framework, along with a reproducible evidence package that reports confidence intervals, convergence traces, and multi-metric comparisons.

Accordingly, this manuscript positions the proposed method as a robust secrecy-aware surrogate for practical IRS-NOMA operation. The comparison against AO is labeled as \emph{literature-inspired} rather than a claim of exact reproduction, while the overall novelty claim rests on the combined formulation, the robust projected solver, and the reproducible evidence stack generated directly from the codebase.
